% Introducción SQA para DIAPCE
% Incluir en el documento maestro con: % Introducción SQA para DIAPCE
% Incluir en el documento maestro con: % Introducción SQA para DIAPCE
% Incluir en el documento maestro con: % Introducción SQA para DIAPCE
% Incluir en el documento maestro con: \input{docs/sqa/introduccion.tex}

\addcontentsline{toc}{section}{Introducción}
El presente plan de aseguramiento de la calidad del software (SQA, por sus siglas en inglés) tiene como propósito garantizar que el sistema propuesto —DIAPCE, una aplicación para el diseño asistido de mezclas de concreto basada en datos experimentales— cumpla con los estándares de calidad necesarios para su correcto funcionamiento, confiabilidad, usabilidad y valor técnico. El SQA se integra como una parte esencial de la gestión del proyecto, asegurando que cada fase, desde el diseño de la base de datos y la carga del dataset hasta la entrega y el mantenimiento, se desarrolle bajo prácticas sistemáticas de verificación y validación.

La disciplina de SQA aporta múltiples ventajas. Permite detectar y corregir errores tempranamente (por ejemplo, inconsistencias en migraciones, validaciones de rangos numéricos como la resistencia objetivo en [24{,}00, 57{,}00] MPa, o flujos en cascada de parámetros ambientales), lo que reduce costos y tiempos de corrección; asegura consistencia en los procesos de desarrollo (estandarizando consultas para predicciones de resistencia a 7/14/28 días y la obtención de códigos de aditivo \textbf{P0}--\textbf{PP3} (\textbf{P0}: sin~aditivo; \textbf{P1}--\textbf{P3}: impermeabilizante 2--4\,\%; \textbf{PP1}--\textbf{PP3}: plastificante 0{,}2--0{,}6\,\%)), generando confianza en los resultados; y mejora la experiencia de los usuarios finales —ingenieros, técnicos de laboratorio y responsables de obra— al ofrecer un producto confiable, trazable y útil para la toma de decisiones. En la estructura del proyecto, el SQA actúa como un eje transversal: acompaña todas las etapas, garantizando que los requisitos técnicos y de dominio (integridad de datos, trazabilidad de decisiones, desempeño de consultas) se cumplan de manera efectiva y medible. Su valor radica en alinear la calidad del producto con las expectativas de los stakeholders, potenciando la adopción y sostenibilidad de la solución.

Este documento resume el enfoque adoptado para asegurar la calidad del software en DIAPCE. Presenta las estrategias, métodos y métricas que se utilizarán para verificar que el sistema cumple con los objetivos planteados: centralizar y depurar el dataset experimental; ofrecer predicciones reproducibles de resistencia; restringir la entrada de datos mediante un flujo de decisión guiado (temperatura → humedad → relación a/c → aditivo); y garantizar proyectos por usuario con trazabilidad y seguridad local. Asimismo, se describe cómo se gestionarán los riesgos (calidad de datos, rendimiento de consultas, compatibilidad multiplataforma), cómo se garantizará la trazabilidad de los requisitos (desde historias de usuario hasta pruebas y artefactos de base de datos) y qué mecanismos de verificación/validación se aplicarán (pruebas unitarias y de integración sobre servicios y consultas, revisiones de migraciones, métricas de cobertura y revisiones de usabilidad).

En síntesis, este plan ofrece una guía clara de cómo se logrará un software confiable, eficiente y alineado con las necesidades del diseño de mezclas de concreto. Está diseñado no solo para el equipo técnico, sino también para que clientes, usuarios y demás interesados comprendan el valor que el aseguramiento de la calidad aporta a DIAPCE y cómo contribuye directamente al éxito de la solución.


\addcontentsline{toc}{section}{Introducción}
El presente plan de aseguramiento de la calidad del software (SQA, por sus siglas en inglés) tiene como propósito garantizar que el sistema propuesto —DIAPCE, una aplicación para el diseño asistido de mezclas de concreto basada en datos experimentales— cumpla con los estándares de calidad necesarios para su correcto funcionamiento, confiabilidad, usabilidad y valor técnico. El SQA se integra como una parte esencial de la gestión del proyecto, asegurando que cada fase, desde el diseño de la base de datos y la carga del dataset hasta la entrega y el mantenimiento, se desarrolle bajo prácticas sistemáticas de verificación y validación.

La disciplina de SQA aporta múltiples ventajas. Permite detectar y corregir errores tempranamente (por ejemplo, inconsistencias en migraciones, validaciones de rangos numéricos como la resistencia objetivo en [24{,}00, 57{,}00] MPa, o flujos en cascada de parámetros ambientales), lo que reduce costos y tiempos de corrección; asegura consistencia en los procesos de desarrollo (estandarizando consultas para predicciones de resistencia a 7/14/28 días y la obtención de códigos de aditivo \textbf{P0}--\textbf{PP3} (\textbf{P0}: sin~aditivo; \textbf{P1}--\textbf{P3}: impermeabilizante 2--4\,\%; \textbf{PP1}--\textbf{PP3}: plastificante 0{,}2--0{,}6\,\%)), generando confianza en los resultados; y mejora la experiencia de los usuarios finales —ingenieros, técnicos de laboratorio y responsables de obra— al ofrecer un producto confiable, trazable y útil para la toma de decisiones. En la estructura del proyecto, el SQA actúa como un eje transversal: acompaña todas las etapas, garantizando que los requisitos técnicos y de dominio (integridad de datos, trazabilidad de decisiones, desempeño de consultas) se cumplan de manera efectiva y medible. Su valor radica en alinear la calidad del producto con las expectativas de los stakeholders, potenciando la adopción y sostenibilidad de la solución.

Este documento resume el enfoque adoptado para asegurar la calidad del software en DIAPCE. Presenta las estrategias, métodos y métricas que se utilizarán para verificar que el sistema cumple con los objetivos planteados: centralizar y depurar el dataset experimental; ofrecer predicciones reproducibles de resistencia; restringir la entrada de datos mediante un flujo de decisión guiado (temperatura → humedad → relación a/c → aditivo); y garantizar proyectos por usuario con trazabilidad y seguridad local. Asimismo, se describe cómo se gestionarán los riesgos (calidad de datos, rendimiento de consultas, compatibilidad multiplataforma), cómo se garantizará la trazabilidad de los requisitos (desde historias de usuario hasta pruebas y artefactos de base de datos) y qué mecanismos de verificación/validación se aplicarán (pruebas unitarias y de integración sobre servicios y consultas, revisiones de migraciones, métricas de cobertura y revisiones de usabilidad).

En síntesis, este plan ofrece una guía clara de cómo se logrará un software confiable, eficiente y alineado con las necesidades del diseño de mezclas de concreto. Está diseñado no solo para el equipo técnico, sino también para que clientes, usuarios y demás interesados comprendan el valor que el aseguramiento de la calidad aporta a DIAPCE y cómo contribuye directamente al éxito de la solución.


\addcontentsline{toc}{section}{Introducción}
El presente plan de aseguramiento de la calidad del software (SQA, por sus siglas en inglés) tiene como propósito garantizar que el sistema propuesto —DIAPCE, una aplicación para el diseño asistido de mezclas de concreto basada en datos experimentales— cumpla con los estándares de calidad necesarios para su correcto funcionamiento, confiabilidad, usabilidad y valor técnico. El SQA se integra como una parte esencial de la gestión del proyecto, asegurando que cada fase, desde el diseño de la base de datos y la carga del dataset hasta la entrega y el mantenimiento, se desarrolle bajo prácticas sistemáticas de verificación y validación.

La disciplina de SQA aporta múltiples ventajas. Permite detectar y corregir errores tempranamente (por ejemplo, inconsistencias en migraciones, validaciones de rangos numéricos como la resistencia objetivo en [24{,}00, 57{,}00] MPa, o flujos en cascada de parámetros ambientales), lo que reduce costos y tiempos de corrección; asegura consistencia en los procesos de desarrollo (estandarizando consultas para predicciones de resistencia a 7/14/28 días y la obtención de códigos de aditivo \textbf{P0}--\textbf{PP3} (\textbf{P0}: sin~aditivo; \textbf{P1}--\textbf{P3}: impermeabilizante 2--4\,\%; \textbf{PP1}--\textbf{PP3}: plastificante 0{,}2--0{,}6\,\%)), generando confianza en los resultados; y mejora la experiencia de los usuarios finales —ingenieros, técnicos de laboratorio y responsables de obra— al ofrecer un producto confiable, trazable y útil para la toma de decisiones. En la estructura del proyecto, el SQA actúa como un eje transversal: acompaña todas las etapas, garantizando que los requisitos técnicos y de dominio (integridad de datos, trazabilidad de decisiones, desempeño de consultas) se cumplan de manera efectiva y medible. Su valor radica en alinear la calidad del producto con las expectativas de los stakeholders, potenciando la adopción y sostenibilidad de la solución.

Este documento resume el enfoque adoptado para asegurar la calidad del software en DIAPCE. Presenta las estrategias, métodos y métricas que se utilizarán para verificar que el sistema cumple con los objetivos planteados: centralizar y depurar el dataset experimental; ofrecer predicciones reproducibles de resistencia; restringir la entrada de datos mediante un flujo de decisión guiado (temperatura → humedad → relación a/c → aditivo); y garantizar proyectos por usuario con trazabilidad y seguridad local. Asimismo, se describe cómo se gestionarán los riesgos (calidad de datos, rendimiento de consultas, compatibilidad multiplataforma), cómo se garantizará la trazabilidad de los requisitos (desde historias de usuario hasta pruebas y artefactos de base de datos) y qué mecanismos de verificación/validación se aplicarán (pruebas unitarias y de integración sobre servicios y consultas, revisiones de migraciones, métricas de cobertura y revisiones de usabilidad).

En síntesis, este plan ofrece una guía clara de cómo se logrará un software confiable, eficiente y alineado con las necesidades del diseño de mezclas de concreto. Está diseñado no solo para el equipo técnico, sino también para que clientes, usuarios y demás interesados comprendan el valor que el aseguramiento de la calidad aporta a DIAPCE y cómo contribuye directamente al éxito de la solución.


\addcontentsline{toc}{section}{Introducción}
El presente plan de aseguramiento de la calidad del software (SQA, por sus siglas en inglés) tiene como propósito garantizar que el sistema propuesto —DIAPCE, una aplicación para el diseño asistido de mezclas de concreto basada en datos experimentales— cumpla con los estándares de calidad necesarios para su correcto funcionamiento, confiabilidad, usabilidad y valor técnico. El SQA se integra como una parte esencial de la gestión del proyecto, asegurando que cada fase, desde el diseño de la base de datos y la carga del dataset hasta la entrega y el mantenimiento, se desarrolle bajo prácticas sistemáticas de verificación y validación.

La disciplina de SQA aporta múltiples ventajas. Permite detectar y corregir errores tempranamente (por ejemplo, inconsistencias en migraciones, validaciones de rangos numéricos como la resistencia objetivo en [24{,}00, 57{,}00] MPa, o flujos en cascada de parámetros ambientales), lo que reduce costos y tiempos de corrección; asegura consistencia en los procesos de desarrollo (estandarizando consultas para predicciones de resistencia a 7/14/28 días y la obtención de códigos de aditivo \textbf{P0}--\textbf{PP3} (\textbf{P0}: sin~aditivo; \textbf{P1}--\textbf{P3}: impermeabilizante 2--4\,\%; \textbf{PP1}--\textbf{PP3}: plastificante 0{,}2--0{,}6\,\%)), generando confianza en los resultados; y mejora la experiencia de los usuarios finales —ingenieros, técnicos de laboratorio y responsables de obra— al ofrecer un producto confiable, trazable y útil para la toma de decisiones. En la estructura del proyecto, el SQA actúa como un eje transversal: acompaña todas las etapas, garantizando que los requisitos técnicos y de dominio (integridad de datos, trazabilidad de decisiones, desempeño de consultas) se cumplan de manera efectiva y medible. Su valor radica en alinear la calidad del producto con las expectativas de los stakeholders, potenciando la adopción y sostenibilidad de la solución.

Este documento resume el enfoque adoptado para asegurar la calidad del software en DIAPCE. Presenta las estrategias, métodos y métricas que se utilizarán para verificar que el sistema cumple con los objetivos planteados: centralizar y depurar el dataset experimental; ofrecer predicciones reproducibles de resistencia; restringir la entrada de datos mediante un flujo de decisión guiado (temperatura → humedad → relación a/c → aditivo); y garantizar proyectos por usuario con trazabilidad y seguridad local. Asimismo, se describe cómo se gestionarán los riesgos (calidad de datos, rendimiento de consultas, compatibilidad multiplataforma), cómo se garantizará la trazabilidad de los requisitos (desde historias de usuario hasta pruebas y artefactos de base de datos) y qué mecanismos de verificación/validación se aplicarán (pruebas unitarias y de integración sobre servicios y consultas, revisiones de migraciones, métricas de cobertura y revisiones de usabilidad).

En síntesis, este plan ofrece una guía clara de cómo se logrará un software confiable, eficiente y alineado con las necesidades del diseño de mezclas de concreto. Está diseñado no solo para el equipo técnico, sino también para que clientes, usuarios y demás interesados comprendan el valor que el aseguramiento de la calidad aporta a DIAPCE y cómo contribuye directamente al éxito de la solución.
